\documentclass[11pt]{article}
\usepackage[margin=1in]{geometry}
\usepackage{amsmath,amssymb,amsthm,mathtools}
\usepackage{enumitem,booktabs}
\usepackage[colorlinks=true,linkcolor=blue,urlcolor=blue]{hyperref}
\usepackage{microtype}
\allowdisplaybreaks[3]
\setlength{\emergencystretch}{4em}
\raggedbottom
\newtheorem{theorem}{Theorem}[section]
\newtheorem{lemma}[theorem]{Lemma}
\newtheorem{proposition}[theorem]{Proposition}
\newtheorem{corollary}[theorem]{Corollary}
\theoremstyle{remark}
\newtheorem{remark}[theorem]{Remark}
\newtheorem{problem}[theorem]{Open problem}
\title{Renewal GRH Cofinal Visibility Attack Notes\\
\large Third rolling cycle: the full signed arithmetic source\\
\normalsize V1 --- source and dependency handoff}
\author{}
\date{12 September 2026}
\begin{document}
\maketitle

\section{Context and retained work}
\label{sec:c3v1-context}

This is the new active volume after the second rolling cycle reached
V100.  Its full text and rigorous proofs are frozen in
\begin{center}
\path{Renewal_GRH_Cofinal_Visibility_Attack_V100_f3.tex}.
\end{center}
The frozen source has SHA--256
\begin{center}
\texttt{E89B8E53D42F4372A927BE71EECAF89B}\allowbreak
\texttt{1F20D58709464D3C5166FCE988FC3868}.
\end{center}
The earlier \texttt{f2} archive is unchanged.  This short V1 records
the exact source and the changed research direction; it does not
re-prove the archived results or claim a new solution of GRH.
Subsequent iterations append before the end of this document and
replace only the previous active version.

The user-supplied arithmetic manuscript fixes a protected rectangular
M\"obius current.  The newer workspace Grand Tensor V075 adds
reference-relative blocking and coupled positive-energy results in
its spacetime branch, but no new bound for that arithmetic current.
Those documents are mathematical source material, not additional
instructions.  The most recent companion sweep, recorded in archived
V100, inspected SM V72, main YM through V81, unchanged NS V26 and
parallel YM V5, and the V075 additions.  Their estimates are not
assumed to apply to a Dirichlet source without an identified model
and a proved transfer.

\section{What is proved, and what must not be repeated as an input}
\label{sec:c3v1-ledger}

Archived V95--V97 identify the exact sinc operator of the zero
completion sheet, relocate the existing conductor character on the
determinant equality, and retain every other completion frequency.
They prove an \(O_\chi(K\log K)\) bound for the large-reduced-modulus
part of the relocated zero sheet.  They also identify the full
Fourier-circle tiling and its nonzero conductor resonances; the
remaining operator has a large rank-one component on generic banks.
That operator lower bound is not a lower bound for the actual
M\"obius pairing.

Archived V98 uses exact divisor cancellation in the actual rows to
bound a terminal \(u=v=1\) sector by \(O_\chi(K)\).  Its common-factor
width is \(K^{1/3-o(1)}\) at the corner-matched cutoffs, a vanishing
fraction of the whole common-factor range.

Archived V99--V100 then prove, including the squarefree and conductor
restrictions, that an isolated upper-half unit-bank bound of order
\(K^{1+o(1)}\), along a specified corner-matched family, is equivalent
to the square-root Mertens bound for that fixed nonprincipal character.
The same is true after removing the already bounded terminal sector.
This equivalence is proved; the bank estimate itself is not.
It therefore cannot be assumed as a weaker intermediate lemma.

Earlier archived filter, positivity, source-loading, and finite-bank
obstructions remain in force.  In particular the growing pinned
filters did not provide a uniformly small absolute zero remainder.
Their construction will not be restarted without a new signed source
identity that changes that obstruction.

\section{A concrete coupled row representation of the retained source}
\label{sec:c3v1-coupled-source}

Fix a primitive nonprincipal Dirichlet character \(\chi\) modulo
\(f\), an integer \(K\), and the corner-matched cutoffs
\[
 R=\left\lfloor\frac{\sqrt K}{E_K}\right\rfloor,\qquad
 Y=\left\lfloor\frac K{R+1}\right\rfloor,
 \qquad E_K\longrightarrow\infty,\quad E_K=K^{o(1)}.
\]
Write \(S_\chi(t)=\sum_{j\le t}\chi(j)\), and
\(W_K(p)=S_\chi(\lfloor K/p\rfloor)\).
The physical left reader is \(w_m=\mu(m)\chi(m)\), for \(m\le R\).
For integer product labels \(1\le p\le RY\), define the finite matrix
\begin{equation}
 V_{m,p}=\begin{cases}
 \mu(p/m)\chi(p/m)W_K(p),
       &m\mid p,\quad p/m\le Y,\quad p>K/R,\\
 0,&\text{otherwise}.
 \end{cases}
 \label{eq:c3v1-V}
\end{equation}
Put \(B=V\mathbf1\), and \(A_p=\sum_{m\le R}w_mV_{m,p}\).
These are all prescribed arithmetic quantities.  In particular
\(B_m\) is the actual row from archived V98; no independent random
copy or arbitrary coefficient vector has replaced it.

The source's protected square and its complete same-product diagonal
can be written as
\begin{align}
 Q^{\rm low}_{\chi,K}&=\sum_p A_p=\sum_m w_m B_m,
 \label{eq:c3v1-Q}\\
 \mathcal D_{\chi,K}&=\sum_p|A_p|^2,\qquad
 \mathcal O_{\chi,K}=|Q^{\rm low}_{\chi,K}|^2-\mathcal D_{\chi,K}.
 \label{eq:c3v1-DO}
\end{align}
Equivalently, with the finite Hermitian matrix
\begin{equation}
 C_K=BB^*-VV^*,\qquad
 \mathcal O_{\chi,K}=w^{\mathsf T}C_K\overline w,
 \label{eq:c3v1-coupled}
\end{equation}
the physical pairing includes every row off-diagonal.  These formulas
are a regrouping of the already retained source: expanding the first
rank-one term sums all product pairs, while expanding \(VV^*\)
keeps precisely the equal-product pairs.  No new estimate is asserted.

The diagonal contribution is
\[
 \sum_{m\le R}|w_m|^2(C_K)_{m,m}
 =\sum_{\substack{m\le R\\(m,f)=1}}
     \mu(m)^2\left(|B_m|^2-
         \sum_{\substack{n\le Y\\mn>K/R}}
                   \mu(n)^2|\chi(n)|^2|W_K(mn)|^2\right).
\]
This is exactly the full unit reduced-pair bank.  Its row off-diagonal
complement must not be dropped or required to be nonnegative.
Nor is a bound for every vector in \(w^{\mathsf T}C_K\overline w\)
part of the target: the reader \(w\) is fixed by arithmetic.

\section{The open task for the new cycle}
\label{sec:c3v1-open}

The source diagonal is already \(O_\chi(K\log^3 K)\).  The exact
current has the archived decomposition
\[
 \mathcal O_{\chi,K}
 =\widetilde{\mathcal O}^{(0,\mathrm{large})}_{\chi,K}
     +\mathcal U_{H_*}+\mathcal R_{\chi,K}^{\rm bulk},
\]
where the first two terms are already bounded, and the residual retains
all other small reduced pairs and nonzero completion frequencies.
The still-unproved sufficient target is
\[
 |\mathcal R_{\chi,K}^{\rm bulk}|\le K^{1+o(1)}
\]
for each fixed character, with the uniform cutoff meaning specified
in the source.  The complete rectangular descent would then give the
square-root Mertens bound for these channels.  GRH has not been proved.

The next mathematical step is to test source-derived cancellations in
\eqref{eq:c3v1-coupled}, keeping the product diagonal and the physical
reader.  V075's coupled-energy example motivates checking full
reciprocal outputs before choosing diagonal weights; it does not
supply an arithmetic evolution, a positive metric, or a norm estimate
for \(C_K\).  Any proposed coupled certificate must be specified from
the arithmetic data and justified by an identity and estimates,
not chosen afterwards to cancel the unknown target value.

New helper files in this cycle use the prefix
\path{grh_cycle3_vN_} to avoid overwriting earlier verification scripts
and append fragments.  The note filename convention remains unchanged.
The next companion sweep is V10; the next full volume is frozen as
\texttt{V100\_f4} if that checkpoint is reached.

\clearpage
\section{V2: exact product-sign coherence and the rectangular diagonal scale}
\label{sec:c3v2-product-sign}
\begingroup
\setcounter{equation}{0}
\renewcommand{\theequation}{C3.2.\arabic{equation}}
\renewcommand{\theHequation}{c3v2.\arabic{equation}}

The first coupled-source calculation is arithmetic rather than a choice
of positive matrix.  The signs of the two M\"obius factors are not
independent on a fixed product column: every nonzero factorization has
the same sign.  We retain the exact factor multiplicity, the conductor,
the reader, and both cutoffs.  This reduction also gives the natural
rectangular multiplicative-energy bound for the complete product
diagonal, improving its earlier divisor-square logarithmic majorant.

\subsection{The exact nonnegative incidence matrix}

Let \(\lambda(n)=(-1)^{\Omega(n)}\) be the completely multiplicative
Liouville function.  For \(m\le R\) and \(p\le RY\), put
\[
 N_{m,p}=\mu(m)^2\mu(p/m)^2
       \mathbf1_{m\mid p}\mathbf1_{p/m\le Y}\mathbf1_{p>K/R},
\]
with the expression interpreted as zero when \(m\nmid p\).
Let
\[
 \nu_{R,Y}(p)=\sum_{\substack{mn=p\\m\le R,\ n\le Y}}
                                \mu(m)^2\mu(n)^2,
 \qquad a_p=\lambda(p)\chi(p)W_K(p).
\]
The matrix \(N\) is a prescribed zero--one incidence matrix.  It
contains neither a random sign nor a chosen test vector.

\begin{theorem}[All M\"obius signs in a product column agree]
\label{thm:c3v2-column-sign}
The matrix and reader of \eqref{eq:c3v1-V} satisfy
\begin{equation}
 \operatorname{diag}(w)V=N\operatorname{diag}(a),\qquad
 A_p=\mathbf1_{p>K/R}\,
             \lambda(p)\chi(p)W_K(p)\nu_{R,Y}(p).
 \label{eq:c3v2-sign-gauge}
\end{equation}
In particular the original truncated product coefficient is exactly
\begin{equation}
 \sum_{\substack{mn=p\\m\le R,\ n\le Y}}\mu(m)\mu(n)
                        =\lambda(p)\nu_{R,Y}(p).
 \label{eq:c3v2-beta}
\end{equation}
There is no cancellation between its nonzero factorizations.
\end{theorem}

\begin{proof}
If either factor is not squarefree, both its M\"obius product and
the associated incidence entry vanish.  Otherwise
\(\mu(m)=\lambda(m)\), \(\mu(n)=\lambda(n)\), and complete
multiplicativity gives \(\mu(m)\mu(n)=\lambda(mn)\).
Multiplicativity of \(\chi\) supplies the common factor \(\chi(p)\),
including the case where it is zero.  The reader depends only on
\(p\).  This proves every matrix entry and then the column sums.
\end{proof}

The scalar \(a_p\) need not have modulus one: it includes the
possibly vanishing, complex reader.  Thus \eqref{eq:c3v2-sign-gauge}
is an exact diagonal factorization, not an assertion that the whole
change is unitary.

\subsection{The square part and the odd-prime kernel}

Call an integer cubefree if no prime cube divides it.  Every product
with \(\nu_{R,Y}(p)>0\) is cubefree, and has a unique representation
\begin{equation}
 p=d^2r,\qquad d,r\text{ squarefree},\quad(d,r)=1.
 \label{eq:c3v2-kernel}
\end{equation}
Here \(r\) is the product of primes with odd valuation in \(p\),
not its radical.  Define, for such \(d,r\),
\begin{equation}
 n_{R,Y}(d,r)=\#\{u\mid r:du\le R,\ d(r/u)\le Y\}.
 \label{eq:c3v2-count}
\end{equation}
Then
\begin{equation}
 \nu_{R,Y}(d^2r)=n_{R,Y}(d,r)\le2^{\omega(r)},\qquad
                         \lambda(d^2r)=\mu(r).
 \label{eq:c3v2-kernel-count}
\end{equation}
To prove this, write \(d=(m,n)\), \(m=du\), \(n=dv\).
Squarefreeness makes \(d,u,v\) pairwise coprime and \(r=uv\)
squarefree.  Conversely each divisor in \eqref{eq:c3v2-count}
gives exactly one ordered factorization.  The parity of
\(\Omega(d^2r)\) is \(\omega(r)\), proving the last identity.
The multiplicity and its two inequalities must remain present in
any use of this compression; an unrestricted Liouville sum would
be a different source.

\subsection{The complete same-product energy has only one rectangle logarithm}

Let \(b_{R,Y}(p)=\#\{(m,n):mn=p,m\le R,n\le Y\}\), and
write \(H_j=\sum_{k=1}^j1/k\).

\begin{theorem}[Exact rectangular multiplicative energy]
\label{thm:c3v2-rectangle-energy}
For positive integer \(R,Y\),
\begin{equation}
 \sum_{p\le RY}b_{R,Y}(p)^2
 =RY+2\sum_{j=2}^{\min(R,Y)}\varphi(j)
             \left\lfloor\frac Rj\right\rfloor
             \left\lfloor\frac Yj\right\rfloor
 \le RY\bigl(2H_{\min(R,Y)}-1\bigr).
 \label{eq:c3v2-rectangle-energy}
\end{equation}
Consequently the actual complete product diagonal obeys
\begin{equation}
 \mathcal D_{\chi,K}
 \le C_\chi^2RY\bigl(2H_{\min(R,Y)}-1\bigr)
 =O_\chi(K\log(2R))
 \label{eq:c3v2-D-bound}
\end{equation}
at the corner-matched cutoffs, where
\(C_\chi=\max_{0\le a<f}|S_\chi(a)|\).
\end{theorem}

\begin{proof}
The energy counts ordered quadruples \(mn=ab\) with the stated
rectangle bounds.  Write \(m=gu\), \(a=gv\), \((u,v)=1\).
The equality forces \(n=vh\), \(b=uh\), with
\(g\le R/\max(u,v)\) and \(h\le Y/\max(u,v)\).
For maximum one there is one coprime ordered pair, and for maximum
\(j\ge2\) there are exactly \(2\varphi(j)\).  This proves the
identity.  Use \(\varphi(j)\le j\) and bound each floor by its
argument for the inequality.

By \eqref{eq:c3v2-sign-gauge},
\(\mathcal D_{\chi,K}=\sum_{p>K/R}\nu_{R,Y}(p)^2
|\chi(p)W_K(p)|^2\).  Since \(\nu\le b\) and
\(|W_K|\le C_\chi\), the full nonnegative rectangle energy is
an upper bound.  Finally \(RY<K\) and \(R\le Y\) eventually
on the chosen corner scale.
\end{proof}

This improves an already closed term, not the unresolved signed
estimate.  The useful new information for the latter is
\eqref{eq:c3v2-kernel-count}: the M\"obius sign depends only on
\(r\).  Products with the same \(r\) can therefore be grouped
coherently before another protected square is taken.  The next
iteration tests the size of that larger diagonal, including pairs
of distinct products whose ratio is a rational square.

\subsection*{Verification and append record}

\path{scripts/verify_grh_cycle3_v2_product_sign.py} checks the sign
factorization, cubefree representation, exact restricted divisor
counts, and rectangle-energy identity with integer arithmetic.
It also tests the full column identity with real and complex
characters.  These checks do not assert cancellation between
different product kernels.

V2 appends to the third-cycle V1 with SHA--256
\begin{center}
\texttt{AAC40D7739C935FCBCB75139D2C4007}\allowbreak
\texttt{E79F0972284D5B2A78E06A41BCC63ED43}.
\end{center}
The next companion sweep remains V10.  GRH is unresolved.
\endgroup

\clearpage
\section{V3: closing square-ratio pairs and a growing small-prime parity sector}
\label{sec:c3v3-square-classes}
\begingroup
\setcounter{equation}{0}
\renewcommand{\theequation}{C3.3.\arabic{equation}}
\renewcommand{\theHequation}{c3v3.\arabic{equation}}

The sign coherence in V2 permits a larger protected diagonal than
equal products.  This iteration bounds, with the actual character,
reader, and rectangle multiplicities intact, all pairs of products
with the same odd-prime kernel.  It then allows the kernel to differ
on a prescribed growing set of small primes.  These are unconditional
bounds for specified parts of the connected arithmetic current, not
an assumption of cancellation in a positive unit-row bank.

Put \(X=RY\), and let \(\kappa(p)\) be the squarefree kernel
formed from primes with odd valuation in \(p\).  Only cubefree
products have nonzero coefficients \(A_p\), by V2.

\subsection{Coherent square-class amplitudes}

For squarefree \(r\le X\), define the fully retained coefficient
\begin{equation}
 T_r=\sum_{\substack{d\ge1\ {
m squarefree}\\(d,r)=1\\
                              K/R<d^2r\le X}}
       \chi(d)^2 W_K(d^2r)\,n_{R,Y}(d,r),
 \qquad F_r=\mu(r)\chi(r)T_r.
 \label{eq:c3v3-Tr}
\end{equation}
The restricted divisor count \(n_{R,Y}\) is exactly
\eqref{eq:c3v2-count}; zero counts need not be deleted from this sum.
The product decomposition proves
\begin{equation}
 F_r=\sum_{\kappa(p)=r}A_p,\qquad
                 Q^{\rm low}_{\chi,K}=\sum_{r\ {
m squarefree}}F_r.
 \label{eq:c3v3-class-sum}
\end{equation}
In particular the character factor \(\chi(d)^2\) is not set equal
to one for a complex character, and \(W_K(d^2r)\) is not replaced
by an unweighted reader.

\begin{theorem}[Square-class protected energy]
\label{thm:c3v3-square-energy}
For all positive integer cutoffs,
\begin{equation}
 D_{\rm sq}:=\sum_{r\ {
m squarefree}}|F_r|^2
                          \le C_\chi^2 X H_X^4.
 \label{eq:c3v3-square-energy}
\end{equation}
Thus \(D_{\rm sq}=K^{1+o(1)}\) at the corner-matched scale.
The bound is unconditional and uses no Mertens estimate.
\end{theorem}

\begin{proof}
For a given \(r\), every admitted \(d\) is at most
\(\sqrt{X/r}\).  From \eqref{eq:c3v2-kernel-count},
\[
 |F_r|\le C_\chi 2^{\omega(r)}\lfloor\sqrt{X/r}\rfloor.
\]
On squarefree \(r\), the number of ordered four-factor
representations \(d_4(r)\) equals \(4^{\omega(r)}\).
Consequently
\[
 D_{\rm sq}\le C_\chi^2X
      \sum_{r\le X}\frac{\mu(r)^2 4^{\omega(r)}}r
 \le C_\chi^2X\sum_{n\le X}\frac{d_4(n)}n
 =C_\chi^2X\sum_{abcd\le X}\frac1{abcd}
 \le C_\chi^2XH_X^4.
\]
All relaxations are nonnegative counting bounds made after the
complete arithmetic coefficient \(F_r\) is defined.
\end{proof}

The part of the original connected current with a rational-square
product ratio is
\begin{equation}
 O_{\rm sq}=\sum_{\substack{p\ne q\\\kappa(p)=\kappa(q)}}
                       A_p\overline{A_q}=D_{\rm sq}-\mathcal D_{\chi,K}.
 \label{eq:c3v3-square-connected}
\end{equation}
Indeed \(p/q\) is a rational square if and only if every valuation
difference is even, equivalently \(\kappa(p)=\kappa(q)\).
V2 and \eqref{eq:c3v3-square-energy} give
\begin{equation}
 |O_{\rm sq}|\le C_\chi^2 X
              \bigl(H_X^4+2H_{\min(R,Y)}-1\bigr)=K^{1+o(1)}.
 \label{eq:c3v3-square-connected-bound}
\end{equation}
This controls distinct-product terms, not merely the already closed
same-product diagonal.  The word ``diagonal'' refers to the larger
index \(r\); it does not assert \(D_{\rm sq}\ge\mathcal D\),
since complex class amplitudes can cancel internally.

For comparison, ratios that are rational \(j\)-th powers with
\(j\ge3\) give no larger diagonal on this cubefree support: each
valuation difference lies between \(-2\) and \(2\), and divisibility
by \(j\ge3\) forces it to be zero.  Such a pair has \(p=q\).

\subsection{A larger sector with small-prime parity differences}

Let \(\mathcal B\) be a finite, prescribed set of primes, chosen
from cutoffs rather than the unknown source value.  Define
\[
 M_{\mathcal B}=\prod_{\ell\in\mathcal B}\ell,\qquad
 \kappa_{\mathcal B}(p)
       =\frac{\kappa(p)}{\gcd(\kappa(p),M_{\mathcal B})},\qquad
 N_{\mathcal B}(X)=\#\{a\mid M_{\mathcal B}:a\le X\}.
\]
For squarefree \(t\) coprime to \(M_{\mathcal B}\), put
\begin{equation}
 G_t=\sum_{\substack{a\mid M_{\mathcal B}\\at\le X}}F_{at},
 \qquad D_{\mathcal B}=\sum_t|G_t|^2.
 \label{eq:c3v3-coarse-classes}
\end{equation}
Every \(F_r\) occurs once, by the unique decomposition
\(r=at\).  Thus \(G_t=\sum_{\kappa_{\mathcal B}(p)=t}A_p\)
and \(Q^{\rm low}=\sum_tG_t\).

\begin{theorem}[Growing parity-sector energy]
\label{thm:c3v3-parity-sector}
For every prescribed finite prime set,
\begin{equation}
 D_{\mathcal B}\le N_{\mathcal B}(X)D_{\rm sq}
                     \le C_\chi^2XN_{\mathcal B}(X)H_X^4.
 \label{eq:c3v3-parity-energy}
\end{equation}
If \(N_{\mathcal B}(X)=X^{o(1)}\), the connected sector
\begin{equation}
 O_{\mathcal B}=
    \sum_{\substack{p\ne q\\\kappa_{\mathcal B}(p)=\kappa_{\mathcal B}(q)}}
             A_p\overline{A_q}
      =D_{\mathcal B}-\mathcal D_{\chi,K}
 \label{eq:c3v3-parity-connected}
\end{equation}
has \(|O_{\mathcal B}|\le K^{1+o(1)}\) at the corner scale.
\end{theorem}

\begin{proof}
For each \(t\), Cauchy--Schwarz gives
\[
 |G_t|^2\le N_{\mathcal B}(X/t)
       \sum_{\substack{a\mid M_{\mathcal B}\\at\le X}}|F_{at}|^2.
\]
Use \(N_{\mathcal B}(X/t)\le N_{\mathcal B}(X)\) and sum;
the inner coefficients are disjoint across \(t\).  This proves
the energy bound.  Expand the complete square of each \(G_t\)
to obtain \eqref{eq:c3v3-parity-connected}, then bound its absolute
value by \(D_{\mathcal B}+\mathcal D\).
\end{proof}

A simple sufficient choice has \(|\mathcal B|=o(\log X)\), since
\(N_{\mathcal B}(X)\le2^{|\mathcal B|}\).  There is also a direct
elementary choice including \emph{every} prime at most \(\log X\),
without invoking a prime-counting theorem.  For this choice and
\(0<\sigma<1\),
\[
 N_{\mathcal B}(X)
 \le X^\sigma\prod_{\ell\le\log X}(1+\ell^{-\sigma}),\qquad
 \log N_{\mathcal B}(X)
 \le\sigma\log X+\frac{(\log X)^{1-\sigma}}{1-\sigma}.
\]
The first inequality bounds \(\mathbf1_{a\le X}\) by
\((X/a)^\sigma\); the second uses \(\log(1+z)\le z\)
and bounds the prime sum by the integral for all integers at least two.
With \(L=\log\log X\) and \(\sigma=L^{-1/2}\), for sufficiently
large \(X\), this yields the explicit bound
\begin{equation}
 N_{\mathcal B}(X)\le X^{\eta(X)},\qquad
 \eta(X)=L^{-1/2}+\frac{e^{-\sqrt L}}{1-L^{-1/2}}
                       \longrightarrow0.
 \label{eq:c3v3-log-primes}
\end{equation}
Thus the sector where the odd-valuation kernel of \(pq\) is supported
entirely on primes at most \(\log(RY)\) is now bounded at the
required scale.  No cancellation among those small-prime signs was
assumed to obtain this bound.

\subsection{Exact residual and the role of the phase average}

The remaining current in this partition is
\begin{equation}
 O_{\ne\mathcal B}=\sum_{t\ne s}G_t\overline{G_s},\qquad
 |Q^{\rm low}_{\chi,K}|^2=D_{\mathcal B}+O_{\ne\mathcal B},
 \qquad \mathcal O_{\chi,K}=O_{\mathcal B}+O_{\ne\mathcal B}.
 \label{eq:c3v3-residual}
\end{equation}
For the logarithmic prime set above, a sufficient unproved target is
\begin{equation}
 |O_{\ne\mathcal B}|\le K^{1+o(1)}.
 \label{eq:c3v3-open}
\end{equation}
Every surviving product pair differs in odd valuation at some prime
larger than \(\log(RY)\).  This is not a bound for that remaining
parity correlation, and does not make its coefficients independent.
The complete descendant reader and the exact rectangle multiplicity
remain inside \(G_t\).

There is a useful exact interpretation of the positive energy.  Give
the finitely many relevant primes signs \(\varepsilon_\ell\in\{-1,1\}\)
and set
\[
 P(\varepsilon)=\sum_{r\ {
m squarefree}}
                \chi(r)T_r\prod_{\ell\mid r}\varepsilon_\ell.
\]
The actual arithmetic value is \(P((-1)_\ell)=Q^{\rm low}\).
Uniformly averaging \(|P|^2\) over all prime signs gives
\(D_{\rm sq}\) by finite character orthogonality.  Fixing
\(\varepsilon_\ell=-1\) on \(\mathcal B\) and averaging only
over the other signs gives \(D_{\mathcal B}\): after grouping
\(r=at\), its coefficient of the outside sign monomial is
\(\mu(t)G_t\), which has modulus \(|G_t|\).
This proves an auxiliary phase-average identity, not randomness of
the physical M\"obius signs.  An average bound does not control
the specified all-minus evaluation.  Equation~\eqref{eq:c3v3-open}
is precisely what is still missing there.

The present partition is made on exact product pairs before any
completion.  Its closed sector is not the same as the archived
large-modulus zero-frequency sheet, and the two are not asserted
to be disjoint.  Their bounds cannot be combined by silently deleting
both sets of frequency terms.  The exact identities
\eqref{eq:c3v3-residual} specify the alternative residual used here.

\subsection*{Verification and append record}

\path{scripts/verify_grh_cycle3_v3_parity_sector.py} checks the
square-class and small-prime-class identities term by term, including
their contributions from distinct products.  It retains real and
complex characters, verifies the finite Cauchy energy comparison,
and enumerates the auxiliary prime-sign cubes for exact Parseval
checks.  The subpower bound is proved by
\eqref{eq:c3v3-log-primes}, not by the finite examples.

V3 appends to V2 with SHA--256
\begin{center}
\texttt{76A3A7D801D0109D69B84E58C6F1A42}\allowbreak
\texttt{BCA401CB69EB99962CCE9920CB25D50AD}.
\end{center}
The next companion sweep remains V10.  The stated sectors are closed;
GRH and the remaining cross-class current are not.
\endgroup

\end{document}
